\chapter{Installazione}
\label{app:installazione}

Questa appendice copre ogni scenario di deployment supportato per FairWindSK,
dalla prima compilazione su un laptop di sviluppo all'avvio automatico kiosk su
un display di plancia dedicato.

\begin{notebox}
Tutti gli scenari presuppongono un server Signal~K in esecuzione da qualche parte
sulla rete di bordo.
Installa e configura Signal~K prima di installare FairWindSK.
\end{notebox}

\section{Prerequisiti}
\label{app:prerequisiti}

\subsection{Prerequisiti comuni per tutte le piattaforme desktop}

\begin{itemize}
  \item CMake 3.16 o successivo
  \item Un compilatore compatibile C++17 supportato dal kit Qt~6 scelto
  \item Git
  \item Accesso a Internet durante la prima compilazione
        (CMake scarica \code{nlohmann/json}, \code{QtZeroConf} e \code{QHotkey})
\end{itemize}

\subsection{Moduli Qt 6 richiesti --- Desktop}

\rowcolors{2}{LtGray}{white}
\begin{longtable}{L{4.5cm}L{9.2cm}}
\toprule
\rowcolor{Navy}
\color{white}\sffamily\bfseries Modulo &
\color{white}\sffamily\bfseries Scopo \\
\midrule
\code{Core}            & Classi base Qt e ciclo eventi. \\
\code{Gui}             & Grafica 2D e gestione font. \\
\code{Widgets}         & Toolkit widget usato in tutta la shell. \\
\code{Concurrent}      & Thread pool per task in background. \\
\code{Network}         & Client HTTP/REST e TLS. \\
\code{WebSockets}      & Stream WebSocket Signal~K. \\
\code{Xml}             & Parsing configurazione e risorse. \\
\code{Svg}             & Rendering icone OpenBridge. \\
\code{Positioning}     & API posizione GPS. \\
\code{WebEngineWidgets}& Browser incorporato per app web Signal~K. \\
\code{VirtualKeyboard} & Tastiera a schermo per deployment touch. \\
\code{PrintSupport}    & Generazione PDF (report diagnostici). \\
\bottomrule
\end{longtable}

\subsection{Moduli Qt 6 richiesti --- Android e iOS}

Le build mobile sostituiscono \code{WebEngineWidgets} con un livello web più leggero:
\code{Qml}, \code{Quick}, \code{QuickWidgets}, \code{WebView}.
Non includono \code{QtZeroConf}, \code{QHotkey}, \code{PrintSupport} o \code{WebEngineWidgets}.

\section{Build Desktop Generica (tutte le piattaforme)}
\label{app:generica}

\begin{lstlisting}
git clone https://github.com/OpenFairWind/FairWindSK.git
cd FairWindSK
cmake -S . -B build
cmake --build build --parallel
cmake --install build
\end{lstlisting}

\begin{tipbox}
Se l'accesso a Internet è limitato durante la build, installa prima
\code{nlohmann\_json} dal gestore pacchetti.
Su Debian/Ubuntu: \code{sudo apt install nlohmann-json3-dev}.
Su macOS con Homebrew: \code{brew install nlohmann-json}.
\end{tipbox}

\section{Installazione su macOS}
\label{app:macos}

\begin{lstlisting}
brew install qt cmake ninja nlohmann-json
cmake -S . -B build -G Ninja \
  -DCMAKE_PREFIX_PATH="$(brew --prefix qt)"
cmake --build build --parallel
open build/FairWindSK.app
\end{lstlisting}

Per la distribuzione fuori dall'albero di build, usa \code{macdeployqt}:

\begin{lstlisting}
macdeployqt build/FairWindSK.app
\end{lstlisting}

\section{Installazione su Linux Desktop}
\label{app:linux}

\begin{lstlisting}
sudo apt update
sudo apt install \
  build-essential cmake ninja-build git \
  nlohmann-json3-dev \
  qt6-base-dev qt6-base-dev-tools \
  qt6-webengine-dev qt6-webengine-dev-tools \
  qt6-websockets-dev qt6-positioning-dev \
  libqt6svg6-dev qt6-virtualkeyboard-dev \
  qt6-base-private-dev libqt6printsupport6 \
  libavahi-compat-libdnssd-dev libnss-mdns avahi-utils

cmake -S . -B build -G Ninja
cmake --build build --parallel
sudo cmake --install build
\end{lstlisting}

\section{Raspberry Pi OS --- Postazione Nautica}
\label{app:rpi-postazione}

\begin{lstlisting}
sudo apt update
sudo apt install \
  build-essential cmake ninja-build git \
  nlohmann-json3-dev \
  qml6-module-qt-labs-folderlistmodel \
  qml6-module-qtquick-window \
  qml6-module-qtquick-layouts \
  qml6-module-qtqml-workerscript \
  libnss-mdns avahi-utils \
  libavahi-compat-libdnssd-dev libxkbcommon-dev \
  qt6-base-dev qt6-base-dev-tools \
  qt6-webengine-dev qt6-webengine-dev-tools \
  qt6-websockets-dev qt6-positioning-dev \
  libqt6svg6-dev \
  qt6-virtualkeyboard-dev \
  libqt6virtualkeyboard6 \
  qt6-virtualkeyboard-plugin

cmake -S . -B build -G Ninja
cmake --build build --parallel
sudo cmake --install build
\end{lstlisting}

\begin{notebox}
Se la build fallisce con \code{Qt6VirtualKeyboardConfigVersionImpl.cmake} mancante,
applica la soluzione nota:
\begin{lstlisting}
sudo touch /usr/lib/aarch64-linux-gnu/cmake/Qt6VirtualKeyboard/\
Qt6VirtualKeyboardConfigVersionImpl.cmake
\end{lstlisting}
Poi riesegui \code{cmake}.
\end{notebox}

\section{Raspberry Pi OS --- Display di Plancia Kiosk}
\label{app:rpi-kiosk}

Questo è lo scenario consigliato per un display di plancia dedicato:
Raspberry~Pi~4 o~5 con touchscreen, FairWindSK che si avvia automaticamente all'avvio
in modalità kiosk.

\subsection{Panoramica}

\begin{enumerate}
  \item Installa FairWindSK come nella Sez.~\ref{app:rpi-postazione}.
  \item Configura la modalità schermo intero in \ui{Impostazioni~> Principale}.
  \item Lo script \code{fairwindsk-startup} gestisce l'avvio automatico.
  \item (Facoltativo) Configura il pannello desktop per l'auto-hide.
  \item Abilita la tastiera virtuale (solo touchscreen).
  \item Riavvia e verifica.
\end{enumerate}

\subsection{Cosa fa fairwindsk-startup}

\begin{enumerate}
  \item Legge \code{fairwindsk.ini} per trovare il percorso di configurazione.
  \item Controlla il flag \code{virtualKeyboard} e imposta \code{QT\_IM\_MODULE}
        di conseguenza.
  \item Attende fino a 45 secondi che l'endpoint del catalogo app Signal~K risponda.
  \item Avvia FairWindSK.
  \item Se FairWindSK esce con codice~1 (pulsante \textit{Riavvia} in
        \ui{Impostazioni~> Sistema}), lo helper lo riavvia immediatamente.
  \item Se FairWindSK esce con codice~0 (pulsante \textit{Esci}),
        lo helper si chiude normalmente.
\end{enumerate}

\subsection{Nascondere il pannello desktop (labwc / Bookworm)}

\begin{lstlisting}
nano ~/.config/wf-panel-pi.ini
\end{lstlisting}

Aggiungi o modifica:

\begin{lstlisting}
[panel]
autohide = true
\end{lstlisting}

Salva il file e riavvia.

\section{Installazione su Windows}
\label{app:windows}

\begin{lstlisting}
git clone https://github.com/OpenFairWind/FairWindSK.git
cd FairWindSK
cmake -S . -B build -G Ninja ^
  -DCMAKE_PREFIX_PATH="C:\Qt\6.x.x\msvc2022_64"
cmake --build build --parallel
cmake --install build
windeployqt build\FairWindSK.exe
build\FairWindSK.exe
\end{lstlisting}

\section{Installazione su Android}
\label{app:android}

Le build Android usano l'implementazione web mobile Qt WebView.
Qt Creator è il flusso di lavoro consigliato.

\begin{enumerate}
  \item Apri Android Studio e installa SDK, NDK e Build-Tools.
  \item In Qt Maintenance Tool, installa i componenti Qt Android.
  \item In Qt Creator, vai a \ui{Preferenze~> Dispositivi~> Android} e verifica
        i percorsi SDK, NDK, JDK e piattaforma.
  \item Clona il repository e apri \code{CMakeLists.txt} in Qt Creator.
  \item Seleziona un kit Android (es.\ \textit{Android Qt~6.x.x Clang arm64-v8a}).
  \item Build con \ui{Build~> Compila progetto}.
  \item Collega un dispositivo Android con modalità sviluppatore abilitata o avvia
        un emulatore.
  \item Esegui con \ui{Build~> Esegui}.
\end{enumerate}

\subsection{Differenze funzionali su Android}

mDNS non disponibile; tasto SHIFT+TAB non disponibile; usa tastiera di sistema Android;
stampa PDF non disponibile; \code{file://} bloccato dal modello a finestra singola.

\section{Installazione su iOS e iPadOS}
\label{app:ios}

Le build iOS usano la stessa implementazione mobile di Android.

\begin{enumerate}
  \item Usa un Mac con Xcode, strumenti da riga di comando Xcode e kit iOS Qt.
  \item Clona il repository e apri \code{CMakeLists.txt} in Qt Creator.
  \item Seleziona un kit iOS.
  \item Build con \ui{Build~> Compila progetto}.
  \item Per il Simulator: scegli un target Simulator e premi Esegui.
  \item Per il dispositivo fisico: collega iPhone/iPad, verifica le impostazioni
        di firma, esegui da Qt Creator.
\end{enumerate}

\section{Configurazione al Primo Avvio}
\label{app:primo-avvio}

\subsection{Posizioni dei file di configurazione}

\rowcolors{2}{LtGray}{white}
\begin{longtable}{L{3.6cm}L{10.1cm}}
\toprule
\rowcolor{Navy}
\color{white}\sffamily\bfseries File &
\color{white}\sffamily\bfseries Contenuto \\
\midrule
\code{fairwindsk.json}  &
  Tutte le impostazioni visibili all'utente: URL connessione, layout barre, catalogo app,
  definizioni widget, preset comfort, unità e mappature percorsi Signal~K. JSON leggibile. \\
\code{fairwindsk.ini}   &
  Impostazioni gestite da Qt: percorso a \code{fairwindsk.json}, flag debug e token
  di autenticazione. Il token è tenuto qui intenzionalmente per non essere mai incluso
  in un'esportazione di configurazione. \\
\bottomrule
\end{longtable}

\subsection{Connessione a Signal K}

\begin{enumerate}
  \item Apri \ui{Impostazioni~> Connessione}.
  \item Inserisci l'indirizzo del server Signal~K,
        ad es.\ \code{http://192.168.1.100:3000}.
  \item Tocca \key{Connetti}. Attendi che lo stato diventi \textit{Live}.
  \item Se il server richiede autenticazione, tocca \key{Richiedi Token} e approva
        la richiesta nell'interfaccia admin Signal~K.
\end{enumerate}

\subsection{Impostazioni essenziali iniziali}

\rowcolors{2}{LtGray}{white}
\begin{longtable}{L{4.0cm}L{9.7cm}}
\toprule
\rowcolor{Navy}
\color{white}\sffamily\bfseries Impostazione &
\color{white}\sffamily\bfseries Dove e cosa impostare \\
\midrule
Unità            & \ui{Impostazioni~> Unità}: nodi (velocità), miglia nautiche (distanza), m o piedi (profondità). \\
Formato coord.   & \ui{Impostazioni~> Principale}: GG\textdegree{}MM.MMM' per uso nautico standard. \\
Modalità finestra & \ui{Impostazioni~> Principale}: Schermo Intero per display dedicato; Finestra o Centrata per postazione. \\
Preset comfort   & \ui{Impostazioni~> Comfort}: Giorno per i test iniziali; cambia in Notte per la guardia. \\
App launcher     & \ui{Impostazioni~> Applicazioni}: organizza plotter, strumenti e altre app nella pagina~1. \\
Tastiera virtuale & \ui{Impostazioni~> Principale} (solo touch): spunta, poi riavvia FairWindSK. \\
\bottomrule
\end{longtable}

\section{Risoluzione Problemi di Installazione}
\label{app:problemi-installazione}

\subsection{FairWindSK non si avvia dopo l'installazione}

\begin{enumerate}
  \item Esegui il binario da un terminale per vedere l'output degli errori:
        \code{FairWindSK} (Linux/macOS) o \code{build\textbackslash FairWindSK.exe} (Windows).
  \item Controlla che le librerie Qt WebEngine siano disponibili.
        Su Linux: \code{ldd build/FairWindSK | grep "not found"}.
  \item Su Raspberry~Pi, controlla la soluzione CMake per la tastiera virtuale
        (Sez.~\ref{app:rpi-postazione}).
\end{enumerate}

\subsection{Nessuna applicazione appare nel launcher}

\begin{enumerate}
  \item Verifica l'indirizzo del server Signal~K in \ui{Impostazioni~> Connessione}.
  \item Conferma che il server esponga \code{/signalk/v1/apps/list} o \code{/skServer/webapps}.
  \item Almeno un'app web deve essere installata sul server.
  \item Abilita il log debug: aggiungi \code{debug=true} a \code{fairwindsk.ini}
        e riavvia FairWindSK.
\end{enumerate}

\subsection{FairWindSK non si avvia automaticamente su Raspberry Pi}

\begin{enumerate}
  \item Verifica che il file di autoavvio sia installato:
        \code{ls /etc/xdg/autostart/fairwindsk-startup.desktop}
  \item Verifica che lo script \code{fairwindsk-startup} sia eseguibile e nel PATH:
        \code{which fairwindsk-startup}
  \item Esegui lo script di avvio manualmente da un terminale per controllare gli errori:
        \code{fairwindsk-startup}
  \item Assicurati che la sessione desktop supporti XDG autostart
        (il desktop predefinito di Raspberry~Pi~OS lo fa).
\end{enumerate}
